\label{sec:dense_method}
We also ran additional validations on LIBERO, comparing terminal-only, phase-only, and combined dense-and-phase rewards. These validations examine complete benchmark instructions. They do not directly measure the success of the separately trained skills in our adaptation pipeline. We will conduct more experiments which show separate skill success rates.

The combined LIBERO reward uses a manually defined sequence of reach, grasp, lift, move, and place phases. For each phase, a progress measure is normalized to the interval from zero to one. Reaching uses gripper--object distance, lifting uses height gained above the initial object position, and transport uses planar object--target distance. Grasping also uses gripper closure and grasp state; placement uses target proximity, release, and environment success. For current progress $q_t$ and the previous maximum $q_{\mathrm{best}}$, the rewarded increment is $\max(q_t-q_{\mathrm{best}},0)$. A completion predicate advances the tracker and adds a phase bonus. The bounded progress increment and phase bonus are combined with terminal feedback, while the native task-success verifier remains unchanged. The phase-only condition sets the continuous-progress weight to zero.

This reward construction requires a progress function, normalization, and relevant state variables for every phase. In our adaptation pipeline, we use dense feedback for the initial simple skills and reuse their binary completion predicates during full-task training.
